% ============================================================================
% Fichier  : partie2/P02_C08_chaine_custody.tex
% Titre    : Chaîne de Custody et Admissibilité
% Auteur   : MINKA MI NGUIDJOÏ Thierry Emmanuel
% Version  : 1.0 -- Juin 2026
% Statut   : RÉDIGÉ
% Sources  : M5_Investigation_Numerique_PNC.docx (§5.4 formulaire, BLOC 2)
%            M3_Cadre_Normatif_National_PNC.docx (art. 52-59 Loi 2010/012)
%            Prologue (affaire MBONGO), Ch. PIU (E4-E6, QG2)
%            ISO/IEC 27037, ACPO 4 principes
% ============================================================================

\chapter{Chaîne de Custody et Admissibilité}
\label{chap:custody}

\epigraph{%
    « La chaîne de custody ne commence pas quand vous ouvrez Autopsy.
    Elle commence quand vous posez les yeux sur la scène de crime. »%
}{--- Principe opérationnel de la forensique judiciaire}

\niveauChapitre{Fondamental}{%
    Aucun prérequis technique. Ce chapitre est essentiellement procédural
    et juridique. Il est le plus important pour l'admissibilité des preuves
    devant les tribunaux camerounais.%
}

\begin{boxobjectifs}
\begin{itemize}
    \item Définir la \termecle{chaîne de custody} et comprendre pourquoi
          elle est la condition fondamentale de l'admissibilité d'une preuve.
    \item Maîtriser le cadre légal camerounais~: articles~52--59 de la
          \loicam{2010/012} et leurs implications procédurales concrètes.
    \item Remplir correctement le \termecle{formulaire de chaîne de custody}
          pour chaque type d'evidence (PC, mobile, serveur, cloud).
    \item Identifier et éviter les cinq erreurs de custody les plus
          fréquentes qui provoquent le rejet des preuves à l'audience.
    \item Appliquer le protocole de custody à chaque étape~:
          collecte, transport, stockage, transfert, analyse, présentation.
    \item Analyser chaque décision de custody à travers le Trilemme~CRO.
\end{itemize}
\end{boxobjectifs}

\bigskip

% ============================================================================
\section*{Le cas qui revient au début}
% ============================================================================

Le Prologue de ce cours s'est ouvert sur l'affaire MBONGO~: un dossier
judiciaire complet --- emails, captures d'écran, historique de navigation ---
anéanti en dix minutes d'audience parce que l'OPJ avait allumé le PC du suspect
avant de l'acquérir. Tout le travail technique était irréprochable. La chaîne
de custody était brisée.

Ce chapitre est la réponse directe à cette affaire. Il ne s'agit pas de
théorie~: chaque procédure décrite ici correspond à une erreur réelle qui
a fait rejeter une preuve numérique devant un tribunal. Lire ce chapitre,
c'est apprendre à ne pas répéter ces erreurs.

\separateur

% ============================================================================
\section{Définition et Fondements de la Chaîne de Custody}
\label{sec:custody-fondements}
% ============================================================================

\subsection{Définition formelle}

\begin{boxdef}
\textbf{Chaîne de custody}\index{chaîne de custody}
(\textit{Chain of custody, chain of evidence}) --- Document chronologique
enregistrant chaque personne ayant eu accès à une preuve, chaque action
effectuée sur cette preuve, et chaque transfert de cette preuve entre
deux personnes, depuis le moment de sa collecte jusqu'à sa présentation
devant la juridiction compétente. La chaîne de custody est la condition
nécessaire de l'\textbf{authenticité} d'une preuve~: elle permet de prouver
que l'objet présenté à l'audience est bien celui qui a été saisi sur la
scène de crime, sans modification.
\end{boxdef}

\subsection{Pourquoi la chaîne de custody prime sur l'analyse technique}

Un rapport d'analyse de 100~pages, produit par les meilleurs outils du
monde, avec des conclusions irréfutables sur le plan technique, perd toute
valeur probatoire si on ne peut pas prouver que l'evidence analysée est
bien celle saisie sur la scène de crime. Voici pourquoi~:

\begin{enumerate}
    \item \textbf{Le tribunal ne voit pas la scène de crime}~: il ne peut
          évaluer les preuves que sur la base des documents qui lui sont
          soumis. Si ces documents ne prouvent pas l'intégrité de la preuve,
          la preuve n'existe pas juridiquement.
    \item \textbf{La défense cherche la rupture}~: un avocat compétent
          examinera systématiquement la chaîne de custody pour trouver
          une période non documentée, un transfert non signé, un hash
          non calculé. Une seule lacune peut suffire.
    \item \textbf{La contamination peut être involontaire}~: dans l'affaire
          MBONGO, l'OPJ n'avait pas l'intention de contaminer les preuves~;
          il ne savait simplement pas que d'allumer le PC entraînerait
          des modifications automatiques. L'ignorance ne protège pas.
\end{enumerate}

\begin{boxCRO}
\textbf{Analyse CRO --- La chaîne de custody}

La chaîne de custody est le mécanisme qui maximise $\mathcal{O}$
(opposabilité) en construisant une preuve que l'evidence n'a pas été
altérée. Elle a un coût~: $\mathcal{R}$ peut être légèrement réduite
(certaines preuves volatiles se perdent pendant le temps de documentation)
et $\mathcal{C}$ est engagée (la documentation révèle qui a eu accès).

Ce coût est non seulement acceptable, il est \emph{constitutif}~: une preuve
dont la custody n'est pas documentée est une preuve qui n'existe pas
juridiquement, quelle que soit sa valeur technique.

\cro{$\downarrow$ assumé}{$\sim$ légère contrainte}{$\uparrow\uparrow$ maximisé}
\end{boxCRO}

% ============================================================================
\section{Cadre Légal Camerounais de la Custody}
\label{sec:custody-legal}
% ============================================================================

\subsection{Articles~52--59 de la Loi~2010/012~: les perquisitions numériques}

La \loicam{2010/012} consacre ses articles~52--59 aux perquisitions et saisies
en matière de cybersécurité. Ces articles définissent le cadre légal de la
collecte de preuves numériques au Cameroun. Leur maîtrise est obligatoire
pour tout OPJ intervenant sur une scène numérique.

\begin{boxjuris}
\textbf{Article~52 --- Autorisation de perquisition~:}
\begin{quote}
Les officiers de police judiciaire peuvent, dans le cadre de la recherche des
infractions \textit{[...]} procéder à la perquisition des systèmes informatiques
et des réseaux de communications électroniques, ainsi qu'à la saisie des données
informatiques qui s'y trouvent, sur autorisation du procureur de la République
ou du juge d'instruction.
\end{quote}
\textbf{Implication opérationnelle~:} toute perquisition numérique exige une
autorisation préalable du procureur ou du juge. Une investigation menée sans
cette autorisation produit des preuves irrecevables, quelles que soient leur
qualité technique et leur pertinence.
\end{boxjuris}

\begin{boxjuris}
\textbf{Article~53, alinéa~1 --- Copies en présence des parties~:}
\begin{quote}
Les perquisitions \textit{[...]} peuvent porter sur des copies réalisées en
présence des personnes qui assistent à la perquisition.
\end{quote}
\textbf{Implication opérationnelle~:} l'acquisition forensique (création de
l'image E01) doit être réalisée en présence du détenteur ou de son représentant.
C'est la traduction légale du rôle DEFR d'ISO~27037~: l'acquisition est un
acte public, transparent, documenté. L'absence des parties pendant l'acquisition
peut suffire à faire contester l'authenticité.
\end{boxjuris}

\begin{boxjuris}
\textbf{Article~55 --- Saisie des données~:}
\begin{quote}
Les données informatiques saisies sont immédiatement copiées et une copie
est remise contre décharge au détenteur. Le support original est restitué
après réalisation de la copie, sauf décision contraire du juge compétent.
\end{quote}
\textbf{Implication opérationnelle~:} en règle générale, l'OPJ acquiert une
copie forensique (image E01) et restitue l'original. Il ne \emph{saisit}
l'original que sur décision judiciaire. Cela signifie que le hash de
l'original doit être calculé \emph{avant} sa restitution et conservé comme
référence d'intégrité.
\end{boxjuris}

\begin{boxjuris}
\textbf{Article~57 --- Réquisitions aux opérateurs~:}
\begin{quote}
Les officiers de police judiciaire peuvent requérir tout opérateur de
communications électroniques, tout fournisseur de services de communication
au public en ligne \textit{[...]} de lui communiquer les données techniques
relatives à l'identification des numéros d'abonnement ou de connexion,
les données relatives au trafic \textit{[...]}.
\end{quote}
\textbf{Implication opérationnelle~:} les logs d'opérateur (MTN, Orange,
Camtel) et les données des plateformes (Meta, Google, WhatsApp) peuvent être
obtenus par réquisition judiciaire. La réponse de l'opérateur constitue une
preuve numérique soumise aux mêmes règles de custody que toute autre evidence.
\end{boxjuris}

\subsection{Tableau de correspondance~: Loi~2010/012 et ISO~27037}

\begin{tableaucours}{p{3.2cm}p{5.5cm}p{4.0cm}}{%
    \tableauHeader{Article Loi~2010} &
    \tableauHeader{Exigence légale} &
    \tableauHeader{Correspondance ISO~27037}
}
Art.~52 & Autorisation judiciaire préalable &
    QG1 du PIU (mandat validé) \\
Art.~53 al.~1 & Copie réalisée en présence des parties &
    DEFR~: acquisition transparente et documentée \\
Art.~55 & Copie immédiate, original restitué & Hash avant restitution~;
    principe 3-2-1 (E5 PIU) \\
Art.~57 & Réquisitions aux opérateurs &
    Collection phase NIST SP~800-86 \\
\end{tableaucours}
\vspace{-0.8em}
\begin{center}{\small\itshape Tableau~8.1 --- Correspondance Loi~2010/012
et standards forensiques internationaux}
\end{center}

% ============================================================================
\section{Le Formulaire de Chaîne de Custody~: Anatomie et Utilisation}
\label{sec:custody-formulaire}
% ============================================================================

\subsection{Principe général~: le formulaire commence au premier contact}

Le formulaire de custody ne se remplit pas après l'acquisition. Il commence
au moment où l'OPJ pose les yeux sur l'evidence --- avant même de la toucher.
La première entrée est la \textbf{description visuelle initiale}, réalisée
sans manipulation. Toutes les entrées suivantes doivent être horodatées et
signées à mesure qu'elles se produisent, pas reconstituées a posteriori.

\subsection{Formulaire type PNC --- version académique complète}

\begin{boxcustody}
\begin{center}
{\bfseries\color{CouleurPrimaire} FORMULAIRE DE CHAÎNE DE CUSTODY}\\[0.5ex]
{\small République du Cameroun --- Police Nationale --- Direction des Investigations}
\end{center}
\vspace{0.5ex}
\begin{tabular}{p{4cm}p{3.5cm}p{4cm}}
\textbf{N° de dossier~:} & \underline{\hspace{3cm}} &
\textbf{N° d'evidence~:} \underline{\hspace{2cm}} \\[0.6ex]
\textbf{Affaire~:} & \multicolumn{2}{l}{\underline{\hspace{7cm}}} \\[0.6ex]
\textbf{Date de collecte~:} & \underline{\hspace{3cm}} &
\textbf{Heure~:} \underline{\hspace{2cm}} \\[0.6ex]
\end{tabular}
\vspace{0.5ex}
\textbf{Description physique de l'evidence~:}\\
Marque~: \underline{\hspace{3.5cm}}
Modèle~: \underline{\hspace{3.5cm}}
N° de série~: \underline{\hspace{3.5cm}}\\
Couleur~: \underline{\hspace{3cm}}
État~: \underline{\hspace{5cm}} \\[0.5ex]
\textbf{Lieu de collecte~:} \underline{\hspace{6cm}}\\[0.5ex]
\textbf{Collecté par~:} Nom~: \underline{\hspace{3cm}}
Grade/matricule~: \underline{\hspace{3cm}}\\[0.5ex]
\textbf{Hash MD5 de l'original~:}\\
\texttt{\underline{\hspace{10cm}}}\\[0.3ex]
\textbf{Hash SHA-256 de l'original~:}\\
\texttt{\underline{\hspace{10cm}}}\\[0.5ex]
\textbf{Hash MD5 de l'image forensique~:}\\
\texttt{\underline{\hspace{10cm}}}\\[0.3ex]
\textbf{Hash SHA-256 de l'image forensique~:}\\
\texttt{\underline{\hspace{10cm}}}\\[0.5ex]
\textbf{Intégrité vérifiée~:}
$\square$~OUI --- les hash correspondent
\quad $\square$~NON --- acquisition à reprendre\\[0.5ex]
\textbf{Outil d'acquisition~:} Nom~: \underline{\hspace{3cm}}
Version~: \underline{\hspace{2cm}}
OS hôte~: \underline{\hspace{2.5cm}}\\[0.8ex]
\textbf{TRANSFERTS DE CUSTODY~:}\\[0.3ex]
\begin{tabular}{p{2.8cm}p{2.8cm}p{2.0cm}p{3.5cm}}
De & À & Date/Heure & Raison \\
\underline{\hspace{2.6cm}} & \underline{\hspace{2.6cm}} &
\underline{\hspace{1.8cm}} & \underline{\hspace{3.3cm}} \\[0.4ex]
\underline{\hspace{2.6cm}} & \underline{\hspace{2.6cm}} &
\underline{\hspace{1.8cm}} & \underline{\hspace{3.3cm}} \\[0.4ex]
\underline{\hspace{2.6cm}} & \underline{\hspace{2.6cm}} &
\underline{\hspace{1.8cm}} & \underline{\hspace{3.3cm}} \\
\end{tabular}
\\[0.8ex]
\textbf{Signature du collecteur~:} \underline{\hspace{5cm}}
\quad Date~: \underline{\hspace{2.5cm}}
\end{boxcustody}

\subsection{Champ par champ~: comment remplir correctement}

\begin{tableaucours}{p{3.5cm}p{8.0cm}}{%
    \tableauHeader{Champ} &
    \tableauHeader{Comment le remplir --- erreurs fréquentes}
}
N° de dossier & Numéro du parquet ou de l'unité~; cohérent avec tous les
    documents du dossier. Ne jamais laisser vide. \\
N° d'evidence & Numérotation séquentielle dans le dossier (EV-001, EV-002~\ldots).
    Ne jamais renuméroter un bien déjà enregistré. \\
Description physique & Suffisamment précise pour qu'un tiers puisse
    identifier l'objet sans l'avoir vu~: marque, modèle, numéro de série,
    couleur, rayures visibles, caches manquants. \\
Hash MD5 + SHA-256 & Calculer les DEUX sur l'original AVANT acquisition~;
    puis les DEUX sur l'image APRÈS acquisition. Les conserver séparément
    de l'image (sur papier et en fichier texte chiffré). \\
Intégrité vérifiée & Cocher OUI uniquement si hash source = hash image.
    Si NON~: stopper, identifier la cause, recommencer l'acquisition. \\
Outil d'acquisition & Nom \emph{exact} + version \emph{exacte}.
    ``FTK Imager'' n'est pas suffisant. ``FTK Imager 4.7.3.81 sur
    Windows~11 Pro~22H2'' est correct. \\
Transferts & Chaque transfert, même temporaire, doit être documenté~:
    remise au laboratoire, transport au greffe, consultation par un avocat,
    retour pour analyse complémentaire. \\
\end{tableaucours}
\vspace{-0.8em}
\begin{center}{\small\itshape Tableau~8.2 --- Guide de remplissage champ par champ}
\end{center}

% ============================================================================
\section{Protocole de Custody par Type d'Evidence}
\label{sec:custody-protocoles}
% ============================================================================

\subsection{PC allumé~: préserver avant de saisir}

\begin{boxPNC}
\textbf{Protocole PNC --- PC allumé (ordre impératif)}

\begin{enumerate}
    \item \textbf{Photographier l'écran} dans son état actuel avant tout geste.
          Horodatage visible (barre des tâches + montre de l'OPJ en main).
    \item \textbf{Capturer la RAM}~: FTK~Imager~$\to$
          \texttt{File~$\to$~Capture Memory}. Support de destination~:
          clé USB distincte, \emph{jamais le disque interne}.
    \item \textbf{Documenter les connexions réseau}~: \texttt{netstat -ano}
          (CMD) ou \texttt{Get-NetTCPConnection} (PowerShell). Photographier.
    \item \textbf{Éteindre proprement} via le menu Windows.
          \emph{Exception}~: si destruction de données en cours
          (chiffrement actif, suppression rapide), débrancher immédiatement
          et documenter le motif.
    \item \textbf{Activer le write blocker} logiciel~:
          \begin{lstlisting}[language=bash_forensique]
reg add HKLM\SYSTEM\CurrentControlSet\Control\StorageDevicePolicies
    /v WriteProtect /t REG_DWORD /d 1 /f
          \end{lstlisting}
    \item \textbf{Acquérir avec FTK~Imager}~: format E01, métadonnées
          complètes (N° dossier, N° evidence, examinateur).
    \item \textbf{Calculer et documenter les hash} MD5 et SHA-256.
          Remplir le formulaire. Sceller l'original.
\end{enumerate}
\end{boxPNC}

\subsection{PC éteint~: ne pas allumer}

\begin{boxPNC}
\textbf{Protocole PNC --- PC éteint}

\begin{enumerate}
    \item \textbf{Ne pas allumer.} Appuyer sur le bouton Power contamine
          la scène avant même que l'OS soit chargé (BIOS, UEFI, Windows
          Write Cache~: tous modifient des données).
    \item Photographier le matériel sous tous les angles, câbles en place.
    \item Activer le write blocker et retirer physiquement le disque dur
          (ou SSD) pour connexion sur le poste d'acquisition.
    \item Acquérir, hasher, documenter. Sceller l'original.
    \item Si le BIOS requiert un démarrage pour identifier le modèle~:
          démarrer sur un Live CD forensique (SIFT, Kali) qui ne monte
          pas les partitions Windows automatiquement.
\end{enumerate}
\end{boxPNC}

\subsection{Téléphone mobile~: mode avion avant tout}

\begin{boxPNC}
\textbf{Protocole PNC --- Téléphone mobile}

\begin{center}\textbf{RÈGLE~1~: MODE AVION EN PREMIER}\end{center}

\begin{enumerate}
    \item \textbf{Téléphone allumé}~: mode avion immédiat (coupe les signaux
          sans effacer les données). Si écran inaccessible~: sac de Faraday
          ou enveloppe d'aluminium ménager hermétique.
    \item \textbf{Ne pas brancher sur un chargeur} avant isolation~: le
          branchement active la synchronisation (certains appareils déclenchent
          une sauvegarde automatique qui modifie des données).
    \item \textbf{Documenter}~: N° IMEI (\texttt{*\#06\#} ou sous la batterie),
          opérateur, état de l'écran, niveau de batterie estimé.
    \item \textbf{Acquisition ADB (Android, code connu)}~:
          \begin{lstlisting}[language=bash_forensique]
# Vérifier connexion et droits
adb devices
# Extraction du système de fichiers (niveau logique)
adb pull /sdcard/ ./extraction/sdcard/
# Liste des applications installées
adb shell pm list packages -f > apps_list.txt
# Sauvegarde complète (si autorisée)
adb backup -all -apk -shared -f backup_mbongo.ab
          \end{lstlisting}
    \item \textbf{Téléphone chiffré, code inconnu}~: conserver éteint,
          transmettre à un laboratoire spécialisé. Ne pas tenter de bruteforce
          sans autorisation judiciaire (art.~52 \loicam{2010/012}).
\end{enumerate}
\end{boxPNC}

\subsection{Données cloud et réponses opérateurs}

Les preuves obtenues par réquisition (art.~57 \loicam{2010/012}) --- logs
d'opérateur MTN, réponses Meta/Google, journaux CloudTrail --- sont
des preuves numériques soumises aux mêmes exigences de custody~:

\begin{enumerate}
    \item \textbf{Documenter la réquisition}~: numéro, date, autorité
          émettrice, périmètre demandé.
    \item \textbf{Documenter la réception}~: date de réception de la
          réponse, nom du fichier, format.
    \item \textbf{Calculer immédiatement le hash} du fichier reçu.
    \item \textbf{Conserver la chaîne documentaire complète}~:
          réquisition + accusé de réception + réponse + hash.
    \item \textbf{Ne pas modifier le fichier reçu}~: travailler sur une
          copie. Le fichier original reste en custody scellée.
\end{enumerate}

% ============================================================================
\section{Les Cinq Erreurs Fatales de Custody}
\label{sec:custody-erreurs}
% ============================================================================

Ces cinq erreurs sont les plus fréquentes dans les dossiers rejetés par les
tribunaux camerounais. Chacune correspond à une règle de custody violée.

\subsection{Erreur 1~: allumer le PC suspect sans acquisition préalable}

\begin{boxalert}
\textbf{L'erreur de l'affaire MBONGO.}

Un OPJ allume l'ordinateur du suspect ``pour constater les faits''. Windows
monte les partitions, crée des fichiers temporaires (\texttt{pagefile.sys},
\texttt{hiberfil.sys}), met à jour le registre, modifie les horodatages d'accès
NTFS. Le hash du disque avant et après allumage diffère. La preuve est
contaminée.

\textbf{Règle absolue~:} sur un PC éteint, ne jamais allumer avant d'avoir
activé le write blocker. Sur un PC allumé, capturer la RAM en premier,
\emph{puis} éteindre proprement, \emph{puis} acquérir le disque.
\end{boxalert}

\subsection{Erreur 2~: travailler sur l'evidence originale}

\begin{boxalert}
\textbf{Analyser directement sur le disque original ou sur l'image sans
write blocker.}

Autopsy modifie-t-il les données en les lisant~? Non, si l'image E01 est
montée en lecture seule. Mais si l'OPJ ouvre l'image avec Windows Explorer,
il crée des fichiers \texttt{Thumbs.db} et met à jour les horodatages d'accès.

\textbf{Règle~:} toute analyse doit être conduite~:
\begin{enumerate}
    \item sur une \emph{copie} de l'image forensique~;
    \item montée en \emph{lecture seule} (Autopsy le fait automatiquement)~;
    \item l'original physique est scellé et inaccessible sauf sur autorisation.
\end{enumerate}
\end{boxalert}

\subsection{Erreur 3~: le transfert non documenté}

\begin{boxalert}
\textbf{``J'ai confié l'ordinateur à mon collègue pendant 2 heures.''}\\
Ce transfert non documenté crée une période de custody inconnue.
L'avocat de la défense peut arguer que l'ordinateur a été manipulé
pendant ces deux heures.

\textbf{Règle~:} tout transfert d'une evidence entre deux personnes, même
temporaire, même entre collègues du même service, est documenté~: date,
heure, raison, signatures des deux parties. Si le transfert a eu lieu sans
documentation~: documenter le fait du transfert a posteriori en indiquant
les circonstances, et re-vérifier l'intégrité par hachage avant et après.
\end{boxalert}

\subsection{Erreur 4~: oublier les versions d'outils}

\begin{boxalert}
\textbf{``Analysé avec Autopsy'' dans le rapport.}

L'avocat demande~: quelle version~? Quel système d'exploitation~? Si
l'analyste ne le sait pas, l'expert de la partie adverse peut arguer que
la version utilisée avait un bug connu qui produisait de faux positifs sur
ce type de fichier.

\textbf{Règle~:} documenter systématiquement dans le formulaire de custody~:
\begin{itemize}
    \item Nom de l'outil (ex.~: FTK Imager)
    \item Version exacte (ex.~: 4.7.3.81)
    \item Système d'exploitation hôte (ex.~: Windows~11 Pro~22H2, build 22621)
    \item Date d'utilisation
\end{itemize}
\end{boxalert}

\subsection{Erreur 5~: les hash non calculés ou non conservés séparément}

\begin{boxalert}
\textbf{``Le hash est dans l'image E01.''}\\
Insuffisant. Si l'image E01 est l'objet de la contestation, le hash qui y
est intégré n'est pas une preuve indépendante.

\textbf{Règle~: les hash doivent être conservés séparément de l'image,
en au moins deux endroits distincts~:}
\begin{enumerate}
    \item Manuscrit sur le formulaire papier de custody (signé et daté)~;
    \item Dans un fichier texte chiffré sur un support distinct~;
    \item Optionnel niveau avancé~: horodatage blockchain ou tiers de confiance.
\end{enumerate}
La règle~: si demain l'image E01 disparaît, les hash doivent subsister
ailleurs pour permettre de certifier une copie de sauvegarde.
\end{boxalert}

% ============================================================================
\section{Custody et Scénario MBONGO~: l'Investigation Correcte}
\label{sec:custody-mbongo}
% ============================================================================

Voici la reconstitution complète et correcte de l'affaire MBONGO du Prologue,
avec le formulaire de custody rempli correctement.

\begin{boxscenario}{Affaire MBONGO --- Formulaire de custody complet}
\textbf{Evidence EV-001 --- PC portable Windows de MBONGO Joseph}

\begin{tabular}{p{4cm}p{3cm}p{4cm}}
\textbf{N° dossier~:} & DOS-2026-034 &
\textbf{N° evidence~:} EV-001 \\
\textbf{Date collecte~:} & 15/03/2026 & \textbf{Heure~:} 14h22 \\
\end{tabular}

\textbf{Description~:} Dell Inspiron 15 3000, S/N GF8T3Y2, noir, rayure 3~cm
sur le bord gauche, autocollant ``SONATRAV'' en bas à gauche.\\

\textbf{Lieu~:} Bureau comptabilité, 3\textsuperscript{e} étage, SONATRAV SA,
Akwa, Douala.\\

\textbf{Collecté par~:} OPJ NKENG Éric, Matricule PNC 4471-B, DCPJ Douala.\\

\textbf{État initial~:} PC allumé, déverrouillé, fenêtre navigateur Edge
ouverte sur portail UBA.\\

\textbf{Actions pré-acquisition~:}
\begin{enumerate}
    \item 14h22 --- Photographies de l'écran (8 photos, horodatées).
    \item 14h24 --- Capture RAM avec FTK~Imager~4.7.3.81~:
          fichier \texttt{MBONGO\_RAM\_20260315.mem} (16~Go).
    \item 14h41 --- \texttt{netstat -ano} documenté (2~connexions ESTABLISHED).
    \item 14h43 --- Arrêt propre Windows.
    \item 14h50 --- Write blocker logiciel activé (registre~: WriteProtect=1).
\end{enumerate}

\textbf{Acquisition~:} FTK~Imager~4.7.3.81, Windows~11 Pro~22H2, format~E01.\\
Durée~: 47~minutes (15h03--15h50). Support destination~: HDD externe 2~To,
S/N WD-WXL1A89K3456.\\

\hash{MD5 original}{a3f7b29c8e1d4f56a2c8b7e3d9f0a142}\\
\hash{SHA-256 original}{3d9c8f2b7e4a1c6d5f8e2b9a7c4d3f1e6b8a2c5d7f9e1b3a4c6d8f0e2b4a6c8}\\
\hash{MD5 image}{a3f7b29c8e1d4f56a2c8b7e3d9f0a142}\\
\hash{SHA-256 image}{3d9c8f2b7e4a1c6d5f8e2b9a7c4d3f1e6b8a2c5d7f9e1b3a4c6d8f0e2b4a6c8}\\

\textbf{Intégrité~:} $\checkmark$~OUI --- hash identiques. Acquisition valide.\\

\textbf{Scellé~:} sachet inviolable n°~DL-2026-0347, 15h52, signature OPJ NKENG.\\

\textbf{Transfert~:} De NKENG~$\to$~Dépôt DCPJ Douala.
Date~: 15/03/2026 17h10. Raison~: stockage sécurisé. Signatures : NKENG +
Responsable dépôt OPJ BIYA Marie.
\end{boxscenario}

\begin{boxscenario}{Affaire MBONGO --- Evidence EV-002 --- Téléphone Android}
\textbf{Evidence EV-002 --- Samsung Galaxy A53 de MBONGO Joseph}

\textbf{Actions~:}
\begin{enumerate}
    \item 14h23 --- Mode avion activé (téléphone déverrouillé, écran actif).
    \item 14h23 --- Photographié~: IMEI~: 357128094771034 visible sur l'écran
          (\texttt{*\#06\#}).
    \item 15h55 --- Extraction ADB~: \texttt{adb pull /sdcard/}
          (\texttt{1~447 fichiers extraits}).
    \item 16h10 --- Hash SHA-256 du répertoire d'extraction~: calculé.
    \item 16h12 --- Scellé~: sachet inviolable n°~DL-2026-0348.
\end{enumerate}

L'ensemble de l'extraction est intégré au dossier avec le formulaire de
custody EV-002 signé par OPJ NKENG.
\end{boxscenario}

% ============================================================================
\section{Custody et Admissibilité~: le Chemin vers le Tribunal}
\label{sec:custody-admissibilite}
% ============================================================================

\subsection{Les quatre conditions d'admissibilité revisitées}

La formule d'admissibilité posée au Chapitre~\ref{chap:fond-theo} prend ici
son sens concret~:

\begin{equation}
\mathrm{Admissible}(P) \iff
\underbrace{\mathrm{Légale}(P)}_{\text{art.~52--59}}
\;\wedge\;
\underbrace{\mathrm{Intègre}(P)}_{\text{hash identiques}}
\;\wedge\;
\underbrace{\mathrm{Authentique}(P)}_{\text{custody continue}}
\;\wedge\;
\underbrace{\mathrm{Pertinente}(P)}_{\text{lien avec les faits}}
\end{equation}

\begin{tableaucours}{p{2.8cm}p{4.5cm}p{5.0cm}}{%
    \tableauHeader{Condition} &
    \tableauHeader{Ce que le tribunal vérifie} &
    \tableauHeader{Document qui le prouve}
}
Légalité & Autorisation judiciaire préalable (art.~52)~;
    absence d'actes coercitifs illicites &
    Réquisition ou mandat judiciaire signé \\
Intégrité & Hash source = hash image~;
    aucune modification après acquisition &
    Formulaire custody + hash conservés séparément \\
Authenticité & Chaîne de custody continue depuis la scène~;
    chaque transfert documenté &
    Formulaire custody complet + signatures \\
Pertinence & Lien démontré entre l'evidence et les faits ~;
    rapport forensique documentant les constatations &
    Rapport forensique (section Constatations) \\
\end{tableaucours}
\vspace{-0.8em}
\begin{center}{\small\itshape Tableau~8.3 --- Conditions d'admissibilité
et documents probatoires correspondants}
\end{center}

\subsection{Ce que l'avocat de la défense cherche --- et comment s'y préparer}

Un avocat de la défense compétent examinera systématiquement~:
\begin{enumerate}
    \item \textbf{L'autorisation judiciaire}~: est-elle antérieure à l'acquisition~?
          Son périmètre couvre-t-il ce qui a été saisi~?
    \item \textbf{Les hash}~: sont-ils calculés sur l'original avant acquisition~?
          Sont-ils identiques entre original et image~?
    \item \textbf{Les périodes non documentées}~: y a-t-il un moment où l'evidence
          n'était pas en custody documentée~? Qui l'avait~? Où~?
    \item \textbf{Les versions d'outils}~: les outils utilisés avaient-ils des
          bugs connus sur ce type de fichier à cette date~?
    \item \textbf{La présence des parties}~: l'acquisition a-t-elle été réalisée
          en présence du détenteur (art.~53 al.~1)~?
\end{enumerate}

Pour chacun de ces points, la réponse est dans le formulaire de custody.
Un formulaire complet et signé est la meilleure défense contre toutes
ces attaques.

\separateur

\begin{boxresume}
\textbf{Ce qu'il faut retenir de ce chapitre~:}
\begin{itemize}
    \item La \textbf{chaîne de custody} est la condition de l'\textbf{authenticité}~:
          elle prouve que l'evidence présentée au tribunal est bien celle saisie
          sur la scène, sans modification.

    \item En droit camerounais~: art.~52 (\textbf{autorisation judiciaire}),
          art.~53 al.~1 (\textbf{acquisition en présence des parties}),
          art.~55 (\textbf{copie + restitution de l'original}),
          art.~57 (\textbf{réquisition aux opérateurs}).

    \item Le \textbf{formulaire de custody} commence au premier contact,
          pas après l'acquisition. Champs obligatoires~: description physique,
          hash MD5+SHA-256 (original ET image), outil+version, tous les transferts.

    \item \textbf{Cinq erreurs fatales}~: allumer le PC sans write blocker~;
          travailler sur l'original~; transfert non documenté~; version d'outil
          non renseignée~; hash non conservés séparément de l'image.

    \item \textbf{Les hash doivent être conservés séparément}~: sur le formulaire
          papier signé ET dans un fichier texte sur support distinct~; jamais
          seulement dans l'image E01.

    \item La chaîne de custody maximise $\mathcal{O}$ (opposabilité) au coût
          d'une légère contrainte sur $\mathcal{C}$ et $\mathcal{R}$~; ce coût
          est non seulement acceptable mais constitutif de l'admissibilité.
\end{itemize}
\end{boxresume}

\bigskip

\begin{boxexo}[title={\ding{45}\quad Exercice~8.1 --- Diagnostic de custody \niveaubase}]
Pour chacune des situations suivantes, identifiez la violation de custody
et son impact sur l'admissibilité~:
\begin{enumerate}[(a)]
    \item Un OPJ saisit un téléphone Android, le branche sur son chargeur
          pendant 20~minutes en attendant les renforts, puis l'emballe.
    \item Le rapport forensique indique~: \textit{``Hash SHA-256~:
          3d9c8f2b...''}. Pas de hash MD5. Pas de hash de l'original.
    \item L'evidence a été transportée du dépôt au laboratoire par un
          OPJ~; aucune entrée dans le formulaire pour ce transport.
    \item L'analyste a ouvert l'image E01 avec Windows Explorer pour
          naviguer rapidement dans les fichiers avant d'ouvrir Autopsy.
\end{enumerate}
\end{boxexo}

\begin{boxexo}[title={\ding{45}\quad Exercice~8.2 --- Remplissage de formulaire \niveaumoyen}]
L'Opération TANGUE (Chapitre~\ref{chap:gr-affaires}) a produit les
éléments suivants lors de l'arrestation de BELLO Armand~:
\begin{itemize}
    \item PC portable HP EliteBook 840, S/N CND7382741, allumé,
          verrouillé par code PIN~;
    \item Clé USB Kingston 16~Go bleue, S/N non visible sur l'objet,
          branchée dans le PC~;
    \item Téléphone Samsung Galaxy S22, allumé, déverrouillé.
\end{itemize}
Rédigez le formulaire de custody pour l'evidence EV-001 (PC HP) en
détaillant~: (a)~les actions pré-acquisition dans l'ordre correct,
(b)~les informations à remplir champ par champ, (c)~le protocole
d'acquisition compte tenu de l'état du PC (allumé, verrouillé).
\end{boxexo}

\begin{boxexo}[title={\ding{45}\quad Exercice~8.3 --- Défense en contre-interrogatoire \niveauexpert}]
Vous êtes l'expert forensique dans l'affaire MVOM. L'avocat de TRANSBOIS SA
(qui cherche à contester votre analyse pour réduire sa responsabilité)
vous pose les questions suivantes~:
\begin{enumerate}[(a)]
    \item \textit{``Comment pouvez-vous prouver que l'image E01 que vous avez
          analysée est bien celle du serveur de TRANSBOIS et non une image
          que vous avez préparée~?''}
    \item \textit{``Le serveur était allumé lors de la saisie. Quelle
          modification du disque a été causée par les actions que vous avez
          menées avant d'acquérir l'image~?''}
    \item \textit{``Vous avez utilisé Volatility~3 pour analyser la RAM.
          Cette version avait-elle des bugs connus en mars~2026 sur les
          systèmes Windows~Server~2022~?''}
\end{enumerate}
Rédigez vos réponses telles que vous les donneriez à l'audience.
Pour chaque réponse, identifiez le document de custody qui vous permet
de l'étayer.
\end{boxexo}

\bigskip

\begin{boxinfo}
\textbf{Transition.} Le Chapitre~\ref{chap:acq} approfondit les
\textbf{outils et techniques d'acquisition}~: FTK~Imager en détail,
formats d'image forensique (E01, DD, AFF), write blockers matériels et
logiciels, acquisition de mémoire RAM, et acquisition mobile avancée.
C'est la mise en \oe uvre technique de ce que ce chapitre a établi
sur le plan procédural.
\end{boxinfo}